\begingroup
\small
\setlength{\tabcolsep}{3pt}
\begin{longtable}{@{}L{0.20\textwidth} L{0.16\textwidth} L{0.16\textwidth} L{0.38\textwidth}@{}}
\caption{Unknowns and OGS output fields used in the reference model.}
\label{tab:unknown-output-fields}\\
\toprule
Field & Symbol & Unit & Mathematical meaning and OGS status\\
\midrule
\endfirsthead
\toprule
Field & Symbol & Unit & Mathematical meaning and OGS status\\
\midrule
\endhead
\bottomrule
\endlastfoot
Displacement &
\(\uu=(u_x,u_y)\) &
\(\mathrm{m}\) &
Primary mechanical unknown and run-local output field.  It describes deformation of
the continuum slice relative to the chosen numerical initial state.\\
Liquid pressure &
\(\pl\) &
\(\mathrm{Pa}\) &
Primary hydraulic unknown and run-local output field.  In the Richards convention
used here, negative liquid pressure means positive suction relative to the
gas/reference pressure.\\
Temperature &
\(\TT\) &
\(\mathrm{K}\) &
Primary thermal unknown in the full TRM process and run-local output field.  In the
active parametrization it is constrained to \(T_0=298.15\,\mathrm{K}\) throughout the
domain and time interval.\\
Liquid saturation &
\(\Sl\) &
-- &
Secondary OGS state/output field computed from pressure through the retention law.
It is the mobile OGS liquid occupancy of the pore volume, not total
hydrogen-bearing water in the clay.\\
\end{longtable}
\endgroup

\begingroup
\small
\setlength{\tabcolsep}{3pt}
\begin{longtable}{@{}L{0.20\textwidth} L{0.16\textwidth} L{0.16\textwidth} L{0.38\textwidth}@{}}
\caption{Material fields and constitutive inputs to the active model.}
\label{tab:material-input-fields}\\
\toprule
Input field or parameter & Symbol & Unit & Material meaning and OGS role\\
\midrule
\endfirsthead
\toprule
Input field or parameter & Symbol & Unit & Material meaning and OGS role\\
\midrule
\endhead
\bottomrule
\endlastfoot
Intrinsic permeability &
\(\KK(\xx)\) &
\(\mathrm{m^2}\) &
Material tensor controlling Darcy mobility before multiplication by relative
permeability and division by viscosity.  The current workflow supplies it through
the run-local mesh field \texttt{k\_i\_rd}.\\
Porosity &
\(\por(\xx)\) &
-- &
Continuum pore-volume fraction used in storage, thermal mixing, and the
model-side water-content proxy.  It is supplied through \texttt{n\_rd} and is fixed
at \(0.105\) in the current field package.\\
Elasticity tensor &
\(\CC\) &
\(\mathrm{Pa}\) &
Orthotropic linear-elastic stiffness used in the mechanical constitutive law.  It is
a fixed constitutive input, not a released field in the present workflow.\\
Biot coefficient and Bishop law &
\(\alpha_B\), \(b(\Sl)\) &
-- &
Fixed pore-pressure coupling inputs.  In the active XML, \(\alpha_B=1\) and
\(b(\Sl)=\Sl\).\\
Liquid properties &
\(\rho_l^0(\pl)\), \(\mu_l\) &
\(\mathrm{kg\,m^{-3}}\), \(\mathrm{Pa\,s}\) &
Liquid density and viscosity used in storage and Darcy mobility.  Under the active
isothermal constraint, density depends on pressure but not on temperature.\\
\end{longtable}
\endgroup

\begingroup
\small
\setlength{\tabcolsep}{3pt}
\begin{longtable}{@{}L{0.20\textwidth} L{0.16\textwidth} L{0.16\textwidth} L{0.38\textwidth}@{}}
\caption{Intermediate and derived quantities used by equations and observation operators.}
\label{tab:intermediate-derived-fields}\\
\toprule
Quantity & Symbol & Unit & Mathematical meaning\\
\midrule
\endfirsthead
\toprule
Quantity & Symbol & Unit & Mathematical meaning\\
\midrule
\endhead
\bottomrule
\endlastfoot
Capillary pressure &
\(\pc=-\pl\) &
\(\mathrm{Pa}\) &
Retention-curve argument under the active sign convention.\\
Relative permeability &
\(k_{\mathrm{rel}}(\Sl)\) &
-- &
Saturation-dependent mobility reduction in the liquid phase.  It is computed from
the fixed relative-permeability law, not fitted independently.\\
Darcy velocity and mass flux &
\(\vl\), \(\ql=\rho_l\vl\) &
\(\mathrm{m\,s^{-1}}\), \(\mathrm{kg\,m^{-2}\,s^{-1}}\) &
Derived hydraulic transport quantities.  They are not pore velocity unless divided
by an explicitly chosen pore-volume factor.\\
Strain and stress &
\(\eps(\uu)\), \(\sig\) &
--, \(\mathrm{Pa}\) &
Mechanical fields derived from displacement, orthotropic elasticity, swelling strain,
and Bishop/Biot pore-pressure coupling.\\
Swelling strain or stress contribution &
\(\eps^{\mathrm{sw}}\), \(p_{\mathrm{sw}}\) &
--, \(\mathrm{Pa}\) &
Saturation-dependent mechanical contribution used inside the active constitutive
law.\\
Volumetric water-content proxy &
\(\theta=\por\Sl\) &
-- &
Model-derived mobile-water proxy.  It needs an observation operator before it can be
compared with NMR, ERT, or Taupe/TDR.\\
\end{longtable}
\endgroup
